\SourceLeaf{113}

\begin{center}\textbf{[WEITERE FRAGMENTE ÜBER DAS PENTAGRAMMA MIRIFICUM.]}\end{center}

\begin{center}\rule{1.5cm}{.4pt}\end{center}

\begin{center}[9.]\end{center}

Die Exponenten der Verjüngung der Hauptaxen der centralen Projectionsellipse sind:

\[
\sqrt{\dfrac{G'-1}{G'}} = \dfrac{2G'-1}{\sqrt{\alpha\beta\gamma\delta\varepsilon}} \quad \text{für die erste Axe } \dfrac{1}{\sqrt{G'}},
\]
\[
\sqrt{\dfrac{G''-1}{G''}} = \dfrac{2G''-1}{\sqrt{\alpha\beta\gamma\delta\varepsilon}} \quad \text{für die zweite Axe } \dfrac{1}{\sqrt{G''}},
\]
oder weil

\[
GG'G'' = -\tfrac14 \alpha\beta\gamma\delta\varepsilon,
\]
\[
(G-1)(G'-1)(G''-1) = -\tfrac14,
\]
\[
(2G-1)(2G'-1)(2G''-1) = -\alpha\beta\gamma\delta\varepsilon
\]
[gilt], die Verjüngung der projicirten Axen $= \dfrac{1}{1-2G}$.

Es ist vortheilhaft, neben den vorigen Grössen $G$, $G'$, $G''$ auch die Wurzeln der Gleichung:

\[
\frac{u^3+u}{uu-1} = \sqrt{\alpha\beta\gamma\delta\varepsilon}
\]
einzuführen. Sind dieselben $\xi$, $\eta$, $\zeta$, so ist:

\[
\xi+\eta+\zeta = -\xi\eta\zeta = \sqrt{\alpha\beta\gamma\delta\varepsilon},
\]
\[
\xi\eta+\eta\zeta+\zeta\xi = 1,
\]
